% A technical recheck of the local splitting step, not a second review report.
% Generalizes the coefficient-ring hypothesis in 04-localization.tex.

\paragraph{Local-cover lemma.}
Let $(R,\mathfrak m)$ be a commutative local ring, and let
$\pi:P\twoheadrightarrow N$ be a surjection from a finite free module
such that $P/\mathfrak mP\to N/\mathfrak mN$ is an isomorphism.
For every surjection $\varphi:F\twoheadrightarrow N$ from a finite free
module there is a decomposition
\[
 F=P\oplus Q,\qquad \varphi=(\pi,0),
\]
with $Q$ finite free.  Indeed lift $\varphi$ to $g:F\to P$.
The map $g$ is surjective modulo $\mathfrak m$, hence surjective by
Nakayama applied to its finite cokernel.  Split $g$ using freeness of $P$.
Its kernel is finite projective and therefore free over $R$.
If $h:P\to F$ is a splitting, then
$\varphi h=\pi gh=\pi$ and $\varphi|_{\ker g}=0$.
Thus this decomposition respects the specified map to $N$, not merely
the isomorphism types of the modules.

\paragraph{Application with all edge maps retained.}
Let $K$ be a finite poset with directed edges compatible with its order.
Suppose a sheaf $\mathcal B$ has finite free stalks, is generated at its
stalks by global sections, is flabby for upper sets, and satisfies
$\mathcal B^e=\mathcal B^y/\alpha_e\mathcal B^y$ for $e:a\to y$,
with the upper map the quotient.
For every $c\in K$, assume the constant quotient model
$\mathcal O_c$ on $\{a:a\leq c\}$ is flabby and generated at stalks.
Assume also that for $a<c$ the intersection
\[
 I_{a,c}=\bigcap_{\substack{e:a\to y\\y\leq c}}\alpha_e R
\]
is proper, and that there is at least one such edge.
All labels are in $\mathfrak m$, so the latter condition implies
$I_{a,c}\subseteq\mathfrak m$.

Process a descending linear extension of $K$.  At vertex $a$ let
\[
 N_a=\operatorname{im}\left(
 \Gamma(K_{>a},\mathcal B)\longrightarrow
 \bigoplus_{e:a\to y}\mathcal B^e\right).
\]
Generation and flabbiness give respectively the two inclusions in
$\operatorname{im}(\varphi_a)=N_a$ for the stalk boundary map
$\varphi_a:\mathcal B^a\to\bigoplus_e\mathcal B^e$.

The already constructed isomorphisms at upper stalks determine the
isomorphisms at every edge out of $a$ by quotienting by its label.
Consequently they determine an isomorphism of boundary images.
For an already present copy of $\mathcal O_c$ with $c>a$, the boundary
image is $R/I_{a,c}$: generation and flabbiness identify it with the
image of its stalk $R$.
Taking their direct sum yields $\pi_a:P_a\twoheadrightarrow N_a$,
which is an isomorphism modulo $\mathfrak m$ since
$I_{a,c}\subseteq\mathfrak m$.

Apply the local-cover lemma to get
$\mathcal B^a=P_a\oplus Q_a$, with boundary map $(\pi_a,0)$.
The first summand extends all existing models and all lower-end maps
out of $a$.  The second supplies $\operatorname{rank}Q_a$ copies of
$\mathcal O_a$.  These vanish at every previously processed vertex.
Edges into $a$ are treated later, using the quotient of this already
fixed stalk isomorphism.  Therefore every edge square, for both
endpoints, commutes at the end of the induction.

\paragraph{What hypotheses were actually used.}
The proof gives
$\mathcal B\simeq\bigoplus_{c\in K}\mathcal O_c^{m_c}$.
Neither dimension two, regularity nor unique factorization was used
in the splitting argument.  It works over any commutative local ring
under the displayed model and boundary assumptions.
In the paper, unique factorization with distinct root primes further
identifies $I_{a,c}$ with the product of those labels.  That stronger
identification is used for kernels of arbitrary edge subsets, not
for the existence of the decomposition.
The assumption that each $\mathcal O_c$ is flabby remains essential to
this argument; it is supplied by the separately checked dihedral
character theorem.  This is an extraction of the paper's method,
not a claim of a new literature theorem.
